\section{Introduction to Groups}

\subsection{Basic Axioms and Examples} \label{sec1.1}

\defi A \textbf{binary operation} on a set $G$ is a function $\star : G \times G \to G$. For $a, b \in G$, we write $a \star b$ instead of $\star(a, b)$.

\defi A binary operation $\star$ on $G$ is \textbf{associative} if $(a \star b) \star c = a \star (b \star c)$ for all $a, b, c \in G$.

\defi Let $\star$ be a binary operation on $G$. Then $a, b \in G$ are said to \textbf{commute} if $a \star b = b \star a$. If every pair of elements in $G$ commute, then $\star$ is said to be \textbf{commutative}.

\defi Let $H \subseteq G$ and $\star$ be a binary operation on $G$. Then $H$ is \textbf{closed} under $\star$ if $a \star b \in H$ for all $a, b \in H$. Note that if $\star$ is associative or commutative on $G$, then it is also associative or commutative on $H$.

\defi A \textbf{group} is an ordered pair $(G, \star)$ with set $G$ and binary operation $\star$ such that the following axioms hold:
\begin{enumerate}
    \item $(a \star b) \star c = a \star (b \star c)$ for all $a, b, c \in G$, i.e., $\star$ is associative,
    \item there exists $e \in G$, called the \textbf{identity} of $G$, such that $e \star a = a \star e = a$ for all $a \in G$, and
    \item for each $a \in G$, there exists $a\inv \in G$ (called the \textbf{inverse} of $a$) such that $a \star a\inv = a\inv \star a = e$.
\end{enumerate}
Moreover, we say $(G, \star)$ is \textbf{abelian} if $\star$ is commutative.

\defi Let $(G, \star)$ and $(H, \diamond)$ be groups. The \textbf{direct product} of $G$ and $H$ is the group $(G \times H, *)$ where the elements are from the Cartesian product and $*$ is defined componentwise:
\[(g_1, h_1) * (g_2, h_2) = (g_1 \star g_2, h_1 \diamond h_2) \quad \text{for all } g_1, g_2 \in G \text{ and } h_1, h_2 \in H.\]

\prop[Group Properties] \label{prop1.1}
Let $(G, \star)$ be a group with $g, h \in G$. Then the following hold:
\begin{enumerate}
    \item The identity of $G$ is unique,
    \item the inverse of $g$ is unique,
    \item $(g\inv)\inv = g$,
    \item $(g \star h)\inv = h\inv \star g\inv$,
    \item for $g_1, g_2, \ldots, g_n \in G$, the expression $g_1 \star g_2 \star \cdots \star g_n$ is unambiguous.
\end{enumerate}
\begin{proofnum}
    \item Let $e$ and $e'$ be two identities of $G$. Then $e = e \star e' = e'$.
    \item Let $h$ and $h'$ be two inverses of $g$. Then $h = h \star e = h \star (g \star h') = (h \star g) \star h' = e \star h' = h'$.
    \item Note that $g \star g\inv = e$ so that $g\inv$ is an inverse of $g\inv$. By the uniqueness of inverses, we have $(g\inv)\inv = g$.
    \item We have $(g \star h) \star (g \star h)\inv = e$. Multiplying on the left by $g\inv$ then $h\inv$ proves the result.
    \item We proceed by induction. Note that $n = 1, 2, 3$ are trivial. Assume for every $k < n$ that $k$ elements $g_1, g_2, \ldots, g_k$ can be unambiguously expressed as $g_1 \star g_2 \star \cdots \star g_k$. Then for $n$ elements, we can split the product $g_1 \star g_2 \star \cdots \star g_n$ into two products of $k$ and $n - k$ elements for some $1 \leq k < n$, and by the inductive hypothesis, these two products are unambiguous. Since $\star$ is associative, the result follows.
\end{proofnum}

\prop[Cancellation Laws] \label{prop1.2}
Let $G$ be a group with $g, h \in G$. Then for any $a, b \in G$, the equations $ag = h$ and $gb = h$ have unique solutions for $a$ and $b$, respectively. In particular,
\begin{enumerate}
    \item if $ag = ah$, then $g = h$, and
    \item if $ga = ha$, then $g = h$.
\end{enumerate}

\begin{proof}
    We have $a = hg\inv$ and $b = g\inv h$, and uniqueness follows from uniqueness of inverses. (i) is satisfied by left-multiplying both sides by $a\inv$, and (ii) is satisfied by right-multiplying both sides by $a\inv$.
\end{proof}

\defi Let $G$ be a group with $g \in G$. The \textbf{order} of $g$, denoted by $|g|$, is the smallest positive integer $n$ such that $g^n = 1$. If no such $n$ exists, then we say $g$ has infinite order.

\defi Let $G = \set{g_1, g_2, \ldots, g_n}$ be a group with $g_1 = 1$. The \textbf{multiplication table}, or \textbf{group table}, of $G$ is the $n \times n$ matrix whose $(i, j)$-entry is $g_i g_j$ for $1 \leq i, j \leq n$.

\newpage

\subsection{Dihedral Groups} \label{sec1.2}

\defi Let $n \geq 3$ be an integer. The \textbf{dihedral group} of order $2n$, denoted by $D_{2n}$, is the group of symmetries of a regular $n$-gon. Two elements of $D_{2n}$ are the rotation $r$ by $2\pi/n$ radians and the reflection $s$ across a line of symmetry.

\note The following are claims about $D_{2n}$:
\begin{enumerate}
    \item $1, r, r^2, \ldots, r^{n - 1}$ are all distinct, and $r^n = 1$ so that $|r| = n$.
    \item $|s| = 2$.
    \item $s \neq r^i$ for any $i$.
    \item $sr^i \neq sr^j$ for all $0 \leq i, j \leq n - 1$ with $i \neq j$, so
    \[D_{2n} = \{1, r, r^2, \ldots, r^{n-1}, s, sr, sr^2, \ldots, sr^{n-1}\}.\]
    \item $rs = sr\inv$.
    \item $r^is = sr^{-i}$ for all $0 \leq i \leq n$.
\end{enumerate}

\begin{proofnum}
    \item Clearly, no two rotations by different angles are the same. Moreover, a $2\pi k$ rotation is the identity if and only if $n \mid k$, hence $|r| = n$.
    \item Clear.
    \item Reflections do not preserve orientation while rotations do, so $s$ cannot be a power of $r$.
    \item If (iv) were true, then it would imply $r^i = r^j$ for some $i \neq j$, which contradicts (i).
    \item We consider two cases:
    \begin{enumerate}
        \item If $n$ is odd, then $s(1) = 1$ and $s(i) = n + 2 - i$.
        \item If $n$ is even, then $s(1) = 1$ and $s(n/2 + 1) = n/2 + 1$. For other $i$, then $s(i) = n + 2 - i$.
    \end{enumerate}
    For $r$, we have $r(n) = 1$ and $r(i) = i + 1$ for $1 \leq i \leq n - 1$. Then $r(s(1)) = r(1) = 2$ and $r(s(2)) = r(n) = 1$. On the other hand, $s(r\inv(1)) = s(n) = 1$ and $s(r\inv(2)) = s(1) = 1$. Hence $rs$ and $sr\inv$ have the same effect on vertices 1 and 2, and they must be the same permutation.
    \item By (v), $i = 1$ holds true. Suppose it is true for $i \in \zp$. Then for $i + 1$, we have $r^{i + 1}s = r(r^is) = r(sr^{-i}) = (sr\inv)r^{-i} = sr^{-(i + 1)}$, and the result follows by induction.
\end{proofnum}

\defi A set of \textbf{generators} of a group $G$ is a subset $S \subseteq G$ such that every element of $G$ can be expressed as a finite product of elements of $S$ and their inverses. We denote this as $G = \gen{S}$. We also say $S$ \textbf{generates} $G$.

\defi Equations that are satisfied by the generators of a group are called \textbf{relations}.

\defi Let $G = \gen S$, and let $R_1, \ldots, R_m$ be equations satisfied by elements of $S \cup \set 1$. Then the generators and relations form a \textbf{presentation} of $G$, denoted by
\[G = \gen{S \mid R_1, \ldots, R_m}.\]
A presentation for $D_{2n}$ is given by
\[D_{2n} = \gen{r, s \mid r^n = s^2 = 1, rs = sr\inv}.\]

\ex The group $X_{2n}$ with presentation
\[X_{2n} = \gen{x, y \mid x^n = y^2 = 1, xy = yx^2}\]
has order at most 6 via the chain of equalities
\[x = xy^2 = (xy)y = (yx^2)y = (yx)(xy) = (yx)(yx^2) = y(xy)x^2 = y(yx^2)x^2 = x^4\]

\ex The group $Y$ with presentation
\[Y = \gen{u, v \mid u^4 = v^3 = 1, uv = v^2u^2}\]
is trivial.

\newpage

\subsection{Symmetric Groups} \label{sec1.3}

\defi Let $\Omega$ be a nonempty set, and let $S_\Omega$ be the set of all permutations of $\Omega$. Then $(S_\Omega, \circ)$ is a group under composition of functions, called the \textbf{symmetric group on} $\Omega$.

\defi Let $\Omega = \set{1, 2, \ldots, n}$. Then $S_\Omega$ is denoted by $S_n$, and we say $S_n$ is the \textbf{symmetric group of degree $n$}.

\defi A \textbf{cycle} is a string of integers that represents a permutation in $S_n$ that cyclically permutes the integers in the cycle and fixes all other integers. The number of integers in a cycle is its \textbf{length}, and if there are $\ell$ integers in the cycle, we call it an \textbf{$\ell$-cycle}. We call two cycles with no common integers \textbf{disjoint cycles}.

\defi The \textbf{cycle decomposition} of some $\sigma \in S_n$ is the product of disjoint cycles that represents $\sigma$. Note that the cycle decomposition of $\sigma$ is unique up to the order of the cycles and the starting point of each cycle.

\defi The \textbf{cycle decomposition algorithm} is a set of steps that can be used to find the cycle decomposition of some $\sigma \in S_n$:

\begin{enumerate}
    \item Start with the smallest integer from $\set{1, 2, \ldots, n}$ that has not appeared in any previous cycle. Typically, we start with $a = 1$ if no other cycle has been formed yet.
    \item Apply $\sigma$ to $a$ to get $\sigma(a)$. If $\sigma(a) = a$, then the cycle is $(a)$, and we can move on to the next smallest integer that has not appeared in any previous cycle. Otherwise, we continue applying $\sigma$ to the result until we get back to $a$, and the cycle is formed by the integers in the order they were generated.
    \item Repeat the above steps until all integers have appeared in some cycle.
    \item Remove all single cycles of the form $(a)$ since they represent the identity permutation on $a$.
\end{enumerate}

\ex Let $n = 13$ and define $\sigma \in S_{13}$ as
    \[\begin{array}{ccccccc}
        \sigma(1) = 12, & \sigma(2) = 13, & \sigma(3) = 3, & \sigma(4) = 1, & \sigma(5) = 11, & \sigma(6) = 9, & \sigma(7) = 5, \\
        \sigma(8) = 10, & \sigma(9) = 6, & \sigma(10) = 4, & \sigma(11) = 7, & \sigma(12) = 8, & \sigma(13) = 2
    \end{array}\]
    \begin{itemize}
        \item Starting the cycle, begin with $a = 1$. Since $\sigma(1) = 12\neq 1$, then $\sigma(12) = 8, \sigma(8) = 10, \sigma(10) = 4$, and $\sigma(4) = 1$. The first cycle is then $(1\ 12\ 8\ 10\ 4)$.
        \item The next smallest number is 2, so $\sigma(2) = 13$ and $\sigma(13) = 2$. Then the second cycle is $(2\ 13)$.
        \item The next smallest number is 3, so $\sigma(3) = 3$, and the third cycle is $(3)$.
        \item The next smallest is 5, so $\sigma(5) = 11, \sigma(11) = 7$, and $\sigma(7) = 5$. Then the fourth cycle is $(5\ 11\ 7)$.
        \item The next smallest is 6, so $\sigma(6) = 9$ and $\sigma(9) = 6$. Then the fifth cycle is $(6\ 9)$.
        \item Remove the single cycle $(3)$.
    \end{itemize}
    Following this, we have $\sigma = (1\ 12\ 8\ 10\ 4)(2\ 13)(5\ 11\ 7)(6\ 9)$.

\note $S_n$ is a non-abelian group for all $n \geq 3$ since the cycles $(1\ 2)$ and $(2\ 3)$ do not commute.

\note Disjoint cycles permute only the integers in the cycle, hence they commute.

\note The order of an $\ell$-cycle is simply $\ell$. The order of a cycle with $k$ disjoint cycles of lengths $\ell_1, \ell_2, \ldots, \ell_k$ is $\lcm(\ell_1, \ell_2, \ldots, \ell_k)$.

\newpage

\subsection{Matrix Groups} \label{sec1.4}

\defi A \textbf{field} is a set $F$ with two binary operations $+$ and $\cdot$ such that $(F, +)$ is abelian with identity 0, and ($F - \set 0, \cdot)$ is abelian with identity 1. Moreover, the following distributive law holds:
\[a \cdot (b + c) = (a \cdot b) + (a \cdot c) \quad \text{for all } a, b, c \in F.\]
Moreover, we denote $F\unt = F - \set 0$.

\defi Let $F$ be a field. The \textbf{general linear group of degree $n$} is the set:
\[\gl_n(F) = \set{A \mid \text{$A$ is an $n \times n$ matrix with entries from $F$ and $\det(A) \neq 0$}}\]
with the binary operation of matrix multiplication.

\note Some recorded results for fields, which are proven later:
\begin{itemize}
    \item if $F$ is a field with $|F| < \infty$, then $|F| = p^k$ for prime $p$ and $k \in \zp$, and
    \item if $|F| = p < \infty$, then $|\gl_n(F)| = (p^n - 1)(p^n - p)(p^n - p^2) \cdots (p^n - p^{n - 1})$.
\end{itemize}

\subsection{The Quaternion Group} \label{sec1.5}

\defi The \textbf{quaternion group}, denoted by $Q_8$, is defined to be the set
\[Q_8 = \set{\pm 1, \pm i, \pm j, \pm k}.\]
It has operation $\cdot$ defined by the following rules:
\begin{enumerate}
    \item $1 \cdot q = q \cdot 1 = q$ for all $q \in Q_8$,
    \item $(-1) \cdot (-1) = 1, (-1) \cdot q = q \cdot (-1) = -q$ for all $q \in Q_8$,
    \item $i^2 = j^2 = k^2 = -1$,
    \item $ij = k, jk = i, ki = j$, and
    \item $ji = -k, kj = -i, ik = -j$.
\end{enumerate}
This is a non-abelian group of order 8.

\subsection{Homomorphisms and Isomorphisms} \label{sec1.6}

\defi Let $(G, \star)$ and $(H, \diamond)$ be groups. A \textbf{homomorphism} is a map $\phi : G \to H$ such that
\[\phi(g \star h) = \phi(g) \diamond \phi(h), \quad \text{for all $g, h \in G$}.\]
$\phi$ is an \textbf{isomorphism} if $\phi$ is a bijection. If there exists an isomorphism from $G$ to $H$, then we say $G$ and $H$ are \textbf{isomorphic}, denoted by $G \cong H$. Moreover, for two isomorphic groups $G$ and $H$, it is to be deduced that
\begin{enumerate}
    \item $|G| = |H|$,
    \item $G$ is abelian if and only if $H$ is abelian, and
    \item for all $g \in G$, then $|g| = |\phi(g)|$.
\end{enumerate}

\ex Let $\Gamma$ and $\Sigma$ be two nonempty sets. Then $S_\Gamma \cong S_\Sigma$ if and only if $|\Gamma| = |\Sigma|$. A sketch is as follows:
\begin{enumerate}
    \item Let $\theta : \Gamma \to \Sigma$ be some bijection. Then the map
    \[\phi : S_\Gamma \to S_\Sigma, \quad \phi(\sigma) = \theta \circ \sigma \circ \theta^{-1}\]
    is an isomorphism.
    \item If $S_\Gamma \cong S_\Sigma$, we note that for finite $\Gamma$ and $\Sigma$ where $|\Gamma| = m$ and $|\Sigma| = n$, then $m! = n!$ implies $m = n$.
\end{enumerate}

\note Let $|G| = n < \infty$ with some presentation, and let $S = \set{s_1, s_2, \ldots, s_k}$ generate $G$. Let $H$ be another group with $\set{r_1, \ldots, r_k} \subseteq H$.
\begin{itemize}
    \item If the $r_i$ satisfy the same relations as the $s_i$, then there is a unique homomorphism $\phi : G \to H$ such that $\phi(s_i) = r_i$ for all $i$.
    \item If $H = \gen{r_1, \ldots, r_k}$ and the $r_i$ satisfy the same relations as the $s_i$, then $\phi$ is surjective.
    \item If $H$ is finite and $|H| = |G|$, then $\phi$ is injective, hence an isomorphism.
\end{itemize}

\newpage

\subsection{Group Actions} \label{sec1.7}

\defi A \textbf{group action} of a group $G$ on a set $A$ is a map $G \times A \to A$, written as $g \cdot a$ for $g \in G$ and $a \in A$, such that
\begin{enumerate}
    \item $g \cdot (h \cdot a) = (gh) \cdot a$ for all $g, h \in G$ and $a \in A$, and
    \item $1 \cdot a = a$ for all $a \in A$.
\end{enumerate}

\defi Let a group $G$ act on a set $A$. For fixed $g \in G$, there exists a permutation $\sigma_g$ on $A$ such that
\[\sigma_g(a) = g \cdot a.\]
The map $g \mapsto \sigma_g$ is a homomorphism from $G$ to $S_A$, called the \textbf{permutation representation} of $G$ on $A$. If instead we have a homomorphism $\phi : G \to S_A$, then we can define a group action of $G$ on $A$ by $g \cdot a = \phi(g)(a)$ for all $g \in G$ and $a \in A$.

\defi Let $G$ act on $A$.
\begin{itemize}
    \item The \textbf{trivial action} of $G$ on $A$ is defined by $ga = a$ for all $g \in G$ and $a \in A$. Moreover, the associated permutation representation is the homomorphism $\phi : G \to S_A$ defined by $\phi(g) = \id$ for all $g \in G$.
    \item If distinct elements of $G$ induce distinct permutations on $A$, then the action is \textbf{faithful}. In this case, the associated permutation representation is injective.
    \item The \textbf{kernel} of the action is the set $\set{g \in G \mid ga = a \text{ for all } a \in A}$.
\end{itemize}

\defi Let $G$ act on itself. Then the \textbf{left regular action} of $G$ on itself is defined by $g \cdot a = ga$ for all $g, a \in G$. Note that this action is faithful, since $ga = ha$ for all $a \in G$ implies $g = h$.